\documentclass[11pt, a4paper]{article}
\usepackage[a4paper, margin=2.5cm]{geometry}
\usepackage{fontspec}
\usepackage[spanish, english, bidi=basic, provide=*]{babel}
\babelprovide[import, onchar=ids fonts]{spanish}
\babelfont{rm}{Noto Sans}
\usepackage{enumitem}
\setlist[itemize]{label=-}
\usepackage{hyperref}

\title{\textbf{Guía de Referencia: Gestión de Identidades y Privilegios en Linux}}
\author{Documentación del Sistema NEXUS-9}
\date{\today}

\begin{document}

\maketitle

\section{Introducción}
La administración de identidades en Linux se basa en la manipulación de archivos estructurales (\texttt{/etc/passwd}, \texttt{/etc/group}, \texttt{/etc/shadow}) mediante herramientas de línea de comandos. Este documento detalla las utilidades estándar POSIX y específicas de GNU/Linux requeridas para auditar y modificar cuentas de usuario y grupos, respetando el principio de mínimo privilegio.

\section{Creación y Modificación de Usuarios}

\subsection{useradd}
\textbf{Descripción:} Crea una nueva cuenta de usuario o actualiza la información predeterminada para los nuevos usuarios. Lee los valores por defecto del archivo \texttt{/etc/default/useradd} y las reglas del sistema de \texttt{/etc/login.defs}.

\textbf{Cuándo ejecutarlo:} Cuando un nuevo operador se incorpora a la red (por ejemplo, al sector-core o a una field-terminal) o cuando un nuevo servicio requiere un usuario dedicado (sin acceso a consola) para ejecutarse de forma segura.

\textbf{Ejemplos prácticos:}
\begin{itemize}
    \item \texttt{sudo useradd -m -s /bin/bash operador1} \\
    El flag \texttt{-m} fuerza la creación del directorio home (\texttt{/home/operador1}), y \texttt{-s} define el intérprete de comandos predeterminado.
    \item \texttt{sudo useradd -r -s /usr/sbin/nologin servicio\_nexus} \\
    Crea una cuenta del sistema (\texttt{-r}) con un UID menor a 1000, sin directorio personal y sin capacidad de iniciar sesión (\texttt{nologin}), ideal para ejecutar demonios y respetar el mínimo privilegio.
\end{itemize}

\subsection{usermod}
\textbf{Descripción:} Modifica los atributos de una cuenta de usuario existente. Permite cambiar el nombre de inicio de sesión, el directorio principal, la fecha de caducidad de la cuenta o su pertenencia a distintos grupos.

\textbf{Cuándo ejecutarlo:} Cuando un operador cambia de rol, necesita nuevos permisos (añadirlo a grupos suplementarios como \texttt{sudo} o \texttt{wheel}) o cuando sus datos de acceso deben suspenderse temporalmente.

\textbf{Ejemplos prácticos:}
\begin{itemize}
    \item \texttt{sudo usermod -aG sudo operador1} \\
    Añade (\texttt{-a} de \textit{append}) al usuario al grupo suplementario \texttt{sudo}. Es \textbf{crítico} usar \texttt{-a} junto a \texttt{-G}; si se omite la \texttt{-a}, el usuario será eliminado silenciosamente de todos los demás grupos a los que pertenecía.
    \item \texttt{sudo usermod -L operador1} \\
    Bloquea (\textit{Lock}) la cuenta del usuario colocando un carácter ``!'' en su contraseña cifrada en \texttt{/etc/shadow}, impidiendo el inicio de sesión sin borrar sus datos.
\end{itemize}

\subsection{userdel}
\textbf{Descripción:} Elimina una cuenta de usuario y los archivos relacionados. Modifica los archivos del sistema para borrar todas las referencias al usuario especificado.

\textbf{Cuándo ejecutarlo:} Cuando un operador abandona definitivamente el proyecto o se desinstala un servicio que requería un usuario dedicado.

\textbf{Ejemplos prácticos:}
\begin{itemize}
    \item \texttt{sudo userdel operador1} \\
    Elimina la cuenta pero deja intacto el directorio \texttt{/home/operador1} y sus archivos.
    \item \texttt{sudo userdel -r operador1} \\
    Elimina al usuario junto con su directorio home (\texttt{-r} de \textit{remove}) y su cola de correo (\textit{mail spool}), realizando una limpieza completa.
\end{itemize}

\section{Gestión de Grupos}

\subsection{groupadd}
\textbf{Descripción:} Crea una nueva cuenta de grupo en el sistema añadiendo la entrada correspondiente en \texttt{/etc/group}.

\textbf{Cuándo ejecutarlo:} Para crear conjuntos lógicos de permisos. Por ejemplo, agrupar a todos los administradores de red o desarrolladores para asignarles permisos de lectura compartidos sobre un directorio específico.

\textbf{Ejemplos prácticos:}
\begin{itemize}
    \item \texttt{sudo groupadd mantenimiento} \\
    Crea un nuevo grupo estándar con el próximo GID disponible.
    \item \texttt{sudo groupadd -r sensores} \\
    Crea un grupo de sistema (\texttt{-r}), asignándole un GID bajo.
\end{itemize}

\subsection{groupmod}
\textbf{Descripción:} Modifica la definición de un grupo en el sistema, alterando su nombre o su GID.

\textbf{Cuándo ejecutarlo:} Cuando las políticas de nomenclatura de la organización cambian o cuando existen conflictos de GID en redes que comparten sistemas de archivos (NFS).

\textbf{Ejemplos prácticos:}
\begin{itemize}
    \item \texttt{sudo groupmod -n ingenieros mantenimiento} \\
    Cambia el nombre del grupo (\texttt{-n} de \textit{new name}) de ``mantenimiento'' a ``ingenieros''.
\end{itemize}

\subsection{groupdel}
\textbf{Descripción:} Elimina un grupo existente. No se puede eliminar el grupo primario de un usuario existente sin eliminar primero al usuario.

\textbf{Cuándo ejecutarlo:} Al finalizar un proyecto específico o reorganizar la jerarquía de roles, limpiando grupos que ya no contienen miembros o no cumplen ninguna función.

\textbf{Ejemplos prácticos:}
\begin{itemize}
    \item \texttt{sudo groupdel ingenieros} \\
    Borra el grupo del sistema. Los archivos que pertenecían a este grupo mantendrán el GID numérico huérfano.
\end{itemize}

\section{Auditoría y Consulta de Identidades}

\subsection{id}
\textbf{Descripción:} Imprime los identificadores de usuario (UID) y de grupo (GID) reales y efectivos, junto con los nombres asociados.

\textbf{Cuándo ejecutarlo:} Constantemente, como herramienta de diagnóstico para verificar ``quién soy'' en la terminal actual, confirmar si el comando \texttt{usermod} tuvo éxito, o validar los permisos efectivos antes de ejecutar tareas críticas.

\textbf{Ejemplos prácticos:}
\begin{itemize}
    \item \texttt{id} \\
    Muestra los detalles del usuario actual.
    \item \texttt{id operador1} \\
    Muestra el UID, el GID primario y todos los grupos suplementarios a los que pertenece ``operador1''.
\end{itemize}

\subsection{groups}
\textbf{Descripción:} Imprime los nombres de los grupos a los que pertenece un usuario. Es una herramienta más simplificada en comparación con \texttt{id}.

\textbf{Cuándo ejecutarlo:} Para obtener rápidamente una lista legible por humanos de las membresías de un usuario.

\textbf{Ejemplos prácticos:}
\begin{itemize}
    \item \texttt{groups operador1} \\
    Salida esperada: \texttt{operador1 : operador1 sudo mantenimiento}
\end{itemize}

\subsection{getent}
\textbf{Descripción:} Obtiene entradas de las bases de datos de servicios de nombres (Name Service Switch), como \texttt{passwd}, \texttt{group}, \texttt{hosts}, entre otras. 

\textbf{Cuándo ejecutarlo:} A diferencia de leer directamente el archivo con \texttt{cat /etc/passwd}, \texttt{getent} consultará todas las bases de datos configuradas, incluyendo usuarios que provienen de servidores de red centralizados (LDAP, Active Directory). Es la forma más fiable de consultar identidades.

\textbf{Ejemplos prácticos:}
\begin{itemize}
    \item \texttt{getent passwd} \\
    Muestra la lista de todos los usuarios reconocidos por el sistema.
    \item \texttt{getent group sudo} \\
    Muestra los detalles del grupo \texttt{sudo}, indicando específicamente quiénes son los usuarios que pertenecen a él en su formato delimitado por comas.
\end{itemize}

\section{Seguridad y Autenticación}

\subsection{passwd}
\textbf{Descripción:} Actualiza los tokens de autenticación (contraseñas) de un usuario. Aplica políticas de complejidad de contraseñas y actualiza el archivo cifrado \texttt{/etc/shadow}.

\textbf{Cuándo ejecutarlo:} Cuando un usuario necesita rotar su credencial de acceso por políticas de seguridad, cuando un administrador asigna una clave inicial, o para forzar la expiración de la misma.

\textbf{Ejemplos prácticos:}
\begin{itemize}
    \item \texttt{passwd} \\
    Permite al usuario actual cambiar su propia contraseña (requiere conocer la clave actual).
    \item \texttt{sudo passwd operador1} \\
    Permite al administrador (\textit{root}) sobrescribir la contraseña de ``operador1'' sin necesidad de conocer la clave anterior.
    \item \texttt{sudo passwd -e operador1} \\
    Hace expirar (\texttt{-e}) la contraseña actual. La próxima vez que ``operador1'' inicie sesión, el sistema lo obligará a configurar una contraseña nueva inmediatamente.
\end{itemize}

\end{document}